\section{A standalone verifier for both finite ranges}\label{app:verifier}
The program below uses only the Python standard library. Its complete
inputs are the two residue classes, the cutoff $10^8$ and the small
summand bounds in Proposition~\ref{prop:finite-certificates}. It constructs
the membership sieve itself. For every target it finds an allowed
complementary index, or raises an exception. The direct test of all
divisors of a small candidate is equivalent to requiring all its prime
factors to be $1$ modulo $3$: products of such primes have every divisor
$1$ modulo $3$.

Membership entries are ordered bytes $0$ or $1$. A witness index is
encoded as a signed little-endian 32-bit integer in increasing target
order. The witness is $i$ in the unrestricted identity and $j$, the index
of $b$, in the mixed identity. The expected SHA256 digests are:

\noindent\textit{Unrestricted membership:}
\begin{quote}\footnotesize\ttfamily
010a6f3752c25b6e58164d309c2173f4b308f5b81180ad96c3e96dfca0994575
\end{quote}
\noindent\textit{Unrestricted least-witness indices:}
\begin{quote}\footnotesize\ttfamily
fff9741cf68c12c795e28aebd94884b2717272dd337de5bd28f07234b658ddf0
\end{quote}
\noindent\textit{Mixed membership:}
\begin{quote}\footnotesize\ttfamily
aadd1a6ebe6f6450f4b7e270d308af6971f1b38d3b388b768c3d28c752518fc0
\end{quote}
\noindent\textit{Mixed least-$b$ indices:}
\begin{quote}\footnotesize\ttfamily
91496589fd9b1be7e9aba7cd77d63cb2c0de4be710564172ded09198c638c49d
\end{quote}
The membership arrays have different lengths because the two target
ranges have different final indices. The digests identify outputs;
correctness rests on the exact sieve and exhaustive witness tests.
The accompanying \href{source.zip}{\texttt{source.zip}} contains this
program as \texttt{verify\_binary\_certificates.py}, the separate NumPy
implementations, and the exact-constant and ramified-count checks.

\begingroup\footnotesize
\verbatiminput{verify_binary_certificates.py}
\endgroup
